\section{AI Compute Density：长等待请求改变 serverless 假设}

传统 web request 追求短延迟：CPU 执行、数据库返回、快速结束。AI 应用常调用外部 model provider，streaming 数十秒甚至数分钟；这段时间本地 CPU 可能空闲。如果仍为每个请求独占完整实例，就会产生巨大浪费。课堂预告的 compute 产品后来演化为 Fluid Compute；当前官方资料强调复用 warm instances、并发处理、active CPU 和更长执行生命周期。

\subsection{CPU-bound 与 I/O-bound 必须分开调度}

\term{CPU-bound} 请求的大部分时间在本地计算，过度并发会争抢 CPU；\term{I/O-bound} 请求的大部分时间等待网络、数据库或模型，适度并发可以在等待间隙复用 CPU。定义请求 $j$ 的 active CPU time 为 $a_j$，wall time 为 $w_j$，则等待比例约为
\[
r_j=1-\frac{a_j}{w_j}.
\]
$r_j$ 高的 workload 更适合安全 multiplex；$r_j$ 低则更需要 scale-out 和资源隔离。

\lecturefigure{13-compute-density.png}{AI workload 要按 CPU activity 和 I/O wait 选择 scale-out 或并发复用。}{本地字幕 24:50--28:10；Vercel 当前 Fluid Compute 官方资料。}

\begin{knowledgebox}{读图：compute density 不等于无上限并发}
I/O-bound 请求仍占用内存、socket、file descriptor、stream buffer 和 event-loop capacity。并发上限必须由 latency SLO、memory、downstream quota 和 noisy-neighbor test 决定。调度器需要观察真实 active CPU 和 wait profile，而不是只按请求数扩容。
\end{knowledgebox}

\begin{warningbox}{课堂预告与当前产品不能混为一谈}
访谈只描述了“根据 CPU-bound / I/O-bound profile 智能提高 compute density”的方向。Fluid Compute 后续加入的 active-CPU billing、warm instance reuse、longer execution 和 background work 是当前官方产品描述，不应倒写成 2025 课堂已经完整实现的细节。
\end{warningbox}

\teachervoice{Rauch 的优先级是性能：不能因为某客户突然流行，就让 noisy neighbor 降低其他客户质量。Compute density 的目标不是单纯压成本，而是在保持隔离和延迟标准下减少 idle waste，并把 capacity tuning 从客户手中移到平台。}

\subsection{本章小结}

AI 把 web workload 从短 CPU burst 推向长等待和 streaming。调度器需要区分 active compute 与 I/O wait，在不破坏隔离的前提下复用 warm capacity。

